% Draft centered on fast photometric and geometric convergence.
% Keep the result sentence below as a placeholder until the cross-sequence
% convergence curves and native-rate evaluation are complete.
\begin{abstract}
3D Gaussian Splatting (3DGS) enables fast rendering, but a fast renderer does
not imply that an incrementally built map becomes accurate quickly. We study
\emph{fast photometric and geometric convergence} in 3DGS SLAM: the rate at
which the appearance and structure of newly observed regions improve under a
sequential input stream and a bounded mapping budget. Two bottlenecks prevent
these objectives from being achieved together. Tracking-oriented keyframes do
not necessarily provide the observations required for rapid appearance
refinement, whereas naively adding every frame dilutes the updates available
to each view. Moreover, photometric loss can be reduced by Gaussians with
incorrect depth or opacity, leaving spurious structure in observed free space.
We address these bottlenecks in a DROID-SLAM-based Gaussian mapping pipeline.
For photometric convergence, we expand the mapping-view set beyond tracking
keyframes at a rate determined by completed mapping updates, and introduce
Entropy-Regularized Count Balancing (ERCB) to correct severe imbalances in
per-view optimization while preserving sampling without replacement. For
geometric convergence, a causal free-space Carve loss uses available depth
evidence to suppress opacity in front of observed surfaces. We evaluate both
forms of convergence through held-out appearance and geometry curves against
elapsed time and completed updates, and separately measure end-to-end online
throughput. This formulation treats real-time operation as the constraint
under which fast map convergence must be achieved, rather than as a substitute
for map quality.
% RESULT SLOT: Add one quantitative sentence only after multi-sequence results
% are frozen. If geometry is preserved but not accelerated, replace "geometric
% convergence" by "geometric consistency" throughout the final paper.
\end{abstract}